\documentclass[eng]{ajceam-class}

\usepackage[utf8]{inputenc}
\usepackage[T2A]{fontenc}
\usepackage[russian,english]{babel}
\usepackage{amsmath,amssymb}
\usepackage{graphicx}
\usepackage{csquotes}
\usepackage{hyperref}

% УБРАТЬ СОВСЕМ:
% \received{dd/mm/yyyy}
% \accepted{dd/mm/yyyy}
% \published{dd/mm/yyyy}
\keywords{RFID, Traceability}

\title{Экономический анализ RFID-трассируемости в цепях поставок скоропортящихся продуктов}
\author{Bicmetov Daniel\\Samara National Research University}
\date{\today}
\publicationdate{\today}

\begin{document}
\maketitle

\begin{abstract}
В статье рассматривается использование автоматизированной системы трассируемости на основе технологии RFID для управления запасами скоропортящихся продуктов в цепи поставок, состоящей из производителя и розничного продавца. [file:5] Особое внимание уделяется влиянию уровня гранулярности трассируемости на затраты, выручку и суммарную прибыль цепи поставок. [file:4] На основе опубликованных моделей и численных экспериментов показывается существование экономически оптимального размера трассируемой партии (Economic Traceability Lot), при котором ожидаемая прибыль в присутствии системы трассируемости превышает прибыль в сценарии без трассируемости. [file:5]
\end{abstract}

\section{Introduction}

Оптимизация логистики скоропортящихся продуктов стала одной из ключевых задач из-за ужесточения требований к безопасности пищевой продукции и обязательной трассируемости в цепях поставок. [file:5] Согласно ISO~9001:2000 и Регламенту ЕС №178/2002 трассируемость определяется как способность проследить историю, применение или местонахождение продукции по зарегистрированным идентификаторам на всех стадиях цепи поставок. [file:4] На практике трассируемость реализуется как сквозной процесс, в котором все участники поддерживают записи о поставщиках и клиентах и обмениваются этой информацией, что позволяет ограничить отзыв только теми партиями, которые реально затронуты дефектом или загрязнением. [file:4]

Пищевые компании рассматривают внедрение систем трассируемости как долгосрочное стратегическое вложение, направленное на повышение доверия потребителей и укрепление репутации бренда, но вместе с тем сталкиваются с существенными затратами на проектирование, внедрение и сопровождение таких систем. [file:5] Дополнительную сложность создают особенности цепей поставок скоропортящихся продуктов, в которых качество сильно зависит от способов сбора урожая, технологических процессов переработки, транспортировки и условий хранения. [file:5]

\section{The costs related to traceability systems}

Себестоимость систем трассируемости включает как разовые инвестиции внедрения, так и текущие эксплуатационные расходы, которые зависят от выбранной технологии идентификации, структуры цепи поставок и объёма фиксируемых данных. [file:4] В литературе принято делить затраты на внедрение и эксплуатацию на пять основных категорий: время и усилия персонала, оборудование и программное обеспечение, обучение, услуги внешних консультантов, расходные материалы, а также сертификация и аудит. [file:4]

Операционные расходы существенно зависят от способа идентификации партий и единиц продукции, а также от необходимости мониторинга параметров качества, таких как температура, влажность и газовый состав. [file:5] Уменьшение размера трассируемой партии повышает объём собираемых данных и число идентификаторов, но одновременно снижает неопределённость при отзывах и позволяет точнее локализовать дефектные товары. [file:4] Для небольших и средних предприятий основной барьер составляет необходимость крупных инвестиций, нехватка квалифицированных кадров и организационная инерция, в то время как крупные компании чаще используют автоматизированные системы с высоким уровнем интеграции. [file:5]

\section{RFID technology for traceability systems}

Технология RFID рассматривается как одна из базовых платформ для построения автоматизированных систем трассируемости в агропродовольственных цепях, поскольку позволяет осуществлять бесконтактную идентификацию без прямой видимости и существенного участия персонала. [file:4] RFID\mbox{-}метки могут быть пассивными, полуактивными и активными, причём пассивные метки используют энергию поля считывателя и имеют практически неограниченный срок службы, а активные снабжены собственной батареей и обеспечивают больший радиус действия и возможность записи измеренных данных. [file:4]

Развитие полупассивных и активных RFID\mbox{-}меток с встраиваемыми датчиками температуры, влажности, освещённости и концентрации газов позволяет непрерывно отслеживать состояние продукции и оценивать её остаточный срок годности с опорой на математические модели деградации качества. [file:5] Для задач трассируемости в цепях поставок обычно используются диапазоны HF и UHF, обеспечивающие компромисс между дальностью чтения, помехоустойчивостью и возможностью работы в сложных средах. [file:5] При этом снижение стоимости пассивных UHF\mbox{-}меток сделало экономически оправданным их применение не только для дорогих продуктов, но и для относительно дешёвых скоропортящихся товаров, хотя цена меток остаётся значимым фактором при массовом внедрении. [file:4]

\section{Supply chain approach and model}

Рассматриваемая модель цепи поставок включает производителя и розничного продавца, между которыми движется общий объём продукции $P$ с долей дефектных единиц $p$, причём продукция сгруппирована в партии размером $n$, что задаёт уровень гранулярности трассируемости. [file:5] Каждая партия снабжается RFID\mbox{-}меткой, содержащей уникальный идентификатор и, при необходимости, параметры, характеризующие степень деградации качества, которые считываются на контрольных точках цепи поставок. [file:5]

На стороне производителя проводится приёмочный тест с фиксированным порогом $th$ по доле дефектных единиц в партии, при превышении которого вся партия отклоняется и не отправляется на основной рынок. [file:5] Партии, не прошедшие тест, перенаправляются на альтернативный рынок по сниженной цене (salvage value), либо на переработку, тогда как партии, прошедшие тест, отгружаются розничному продавцу, где оставшийся дефект приводит к затратам на утилизацию. [file:5] Вероятность отклонения партии $R$ формализуется через распределение числа дефектных единиц в партии и задаётся как
\begin{equation}
R = P\{ X \geq n \cdot th \} = 1 - P\{ X < n \cdot th \},
\end{equation}
где $X$ --- случайная величина, описывающая количество дефектных единиц в партии, обычно моделируемая биномиальным распределением с параметрами $n$ и $p$. [file:5]

\subsection{Economic Traceability Lot formulation}

Для оценки экономической эффективности системы трассируемости вводится набор параметров: удельная стоимость утилизации $C_d$, удельная стоимость транспортировки $C_{tr}$, стоимость RFID\mbox{-}метки $C_{tag}$, salvage\mbox{-}стоимость $S_v$ и цена реализации единицы продукции $P_r$. [file:5] Ожидаемое количество отклонённых продуктов, а также ожидаемое количество реализованных качественных и некачественных единиц выражаются через вероятность $R$, размер партии $n$ и исходную долю дефектных единиц $p$. [file:5]

На основе этих величин формируются выражения для ожидаемых совокупных затрат на транспортировку, утилизацию, метки и ожидаемой выручки в сценарии с трассируемостью. [file:5] Итоговая целевая функция представляет собой ожидаемую прибыль цепи поставок как функцию размера партии $n$ и порога $th$, при этом поиск максимума этой функции позволяет определить экономически оптимальный размер трассируемой партии, обозначаемый как Economic Traceability Lot. [file:5]

\section{Numerical analysis and discussion}

Численный анализ модели показывает, что при фиксированном пороге $th$ и нулевой стоимости метки $C_{tag}=0$ ожидаемая прибыль в присутствии трассируемости возрастает с увеличением размера партии $n$ до определённого значения, после чего начинает снижаться и асимптотически приближается к прибыли в сценарии без трассируемости. [file:5] Это поведение указывает на существование оптимального размера партии, при котором достигается максимум прибыли и одновременно обеспечивается требуемый уровень контроля качества. [file:5]

В более общем случае, когда стоимость метки отлична от нуля и доля дефектных единиц $p$ может быть сопоставима с порогом $th$, аналитическое нахождение максимума затруднено, поэтому используются графические и численные методы оптимизации. [file:5] Проведённый чувствительный анализ показывает, что при росте стоимости утилизации и увеличении доли дефекта целесообразно уменьшать размер партии, переходя к более детальной трассируемости, тогда как при низкой стоимости брака и высокой стоимости меток выгоднее использовать более крупные партии. [file:5]

\section{Conclusions}

Представленная модель демонстрирует, что внедрение системы трассируемости на основе RFID в цепях поставок скоропортящихся продуктов может одновременно выполнять функцию инструмента повышения безопасности пищевой продукции и механизма оптимизации экономических показателей цепи. [file:4] Наличие экономически оптимального размера трассируемой партии позволяет обосновать выбор уровня гранулярности, исходя из реальных параметров затрат, качества и регуляторных требований. [file:5]

Дальнейшее развитие подхода связано с учётом ошибок контроля (ложное принятие и ложное отклонение партии), а также с включением фиксированных инфраструктурных затрат и расширенных топологий цепей поставок. [file:5] Такие расширения помогут более точно оценивать инвестиционные решения в области цифровизации логистики и выбирать архитектуры систем ``умной'' трассируемости, адаптированные к специфике конкретных продуктов и рынков. [file:4]

\bibliographystyle{IEEEtran}
\begin{thebibliography}{99}

\bibitem{DabbeneGay2011}
F.~Dabbene and P.~Gay,
``Food traceability systems: Performance evaluation and optimization,''
\emph{Computers and Electronics in Agriculture}, vol.~75, pp.~139--146, 2011. [file:4]

\bibitem{Gandino2009}
F.~Gandino, B.~Montrucchio, M.~Rebaudengo, and E.~R.~Sanchez,
``On improving automation by integrating RFID in the traceability management of the agri-food sector,''
\emph{IEEE Transactions on Industrial Electronics}, vol.~56, no.~7, pp.~2357--2365, 2009. [file:4]

\bibitem{Karlsen2012}
K.~M.~Karlsen, B.~Dreyer, P.~Olsen, and E.~O.~Elvevoll,
``Granularity and its role in implementation of seafood traceability,''
\emph{Journal of Food Engineering}, vol.~112, pp.~78--85, 2012. [file:4]

\bibitem{Pouliot2008}
S.~Pouliot and D.~A.~Sumner,
``Traceability, liability, and incentives for food safety and quality,''
\emph{American Journal of Agricultural Economics}, vol.~90, no.~1, pp.~15--27, 2008. [file:5]

\end{thebibliography}

\end{document}